% Part: many-valued-logic
% Chapter: sequent-calculus
% Section: introduction

\documentclass[../../../include/open-logic-section]{subfiles}

\begin{document}

\olfileid[id]{mvl}{seq}{int}

\olsection{Pengantar}

Kalkulus sekuen bagi logika klasik merupakan sistem !!{derivation} yang
efisien dan sederhana. Jika suatu logika bernilai banyak didefinisikan oleh
matriks dengan nilai kebenaran berhingga---yakni, jika $V$ berhingga---kita
dapat memberikan suatu kalkulus sekuen bagi logika tersebut. Gagasan mengenai
cara melakukannya berasal dari pertimbangan atas makna sekuen dan bentuk aturan
inferensi dalam kasus klasik.

Ingat kembali bahwa suatu sekuen
\begin{align*}
    !A_1, \dots, !A_m & \Sequent !B_1, \dots, !B_n
\intertext{dapat ditafsirkan sebagai !!{formula}}
    (!A_1 \land \cdots \land !A_m) & \lif (!B_1 \lor \cdots \lor
!B_n).
\end{align*}
Dengan kata lain, !!^a{valuation}~$\pAssign{v}$ \emph{memenuhi} suatu sekuen
$\Gamma\Sequent\Delta$ jika dan hanya jika $\pValue{v}(!A)=\False$ bagi
suatu $!A\in\Gamma$ atau $\pValue{v}(!A)=\True$ bagi suatu
$!A\in\Delta$. Dalam penafsiran ini, sekuen awal $!A\Sequent !A$ selalu
dipenuhi, sebab $\pValue v(!A)=\True$ atau $\pValue v(!A)=\False$.

Berikut adalah aturan inferensi bagi implikasi dalam~$\Log{LK}$, dengan
formula samping $\Gamma$, $\Delta$ dihilangkan:

\begin{defish}
    \Axiom$ \fCenter !A$
    \Axiom$ !B \fCenter $
    \RightLabel{\LeftR{\lif}}
    \BinaryInf$ !A \lif !B \fCenter $
    \DisplayProof
    \hfill
    \Axiom$ !A \fCenter !B$
    \RightLabel{\RightR{\lif}}
    \UnaryInf$ \fCenter !A \lif !B $
    \DisplayProof
\end{defish}

Jika kita menerapkan penafsiran semantik sekuen di atas, aturan
$\LeftR{\lif}$ dapat dibaca sebagai pernyataan bahwa jika
$\pValue v(!A)=\True$ dan $\pValue v(!B)=\False$, maka
$\pValue v(!A\lif !B)=\False$. Serupa dengan itu, aturan
$\RightR{\lif}$ mengatakan bahwa jika $\pValue v(!A)=\False$ atau
$\pValue v(!B)=\True$, maka $\pValue v(!A\lif !B)=\True$. Bahkan,
implikasi-implikasi tersebut sesungguhnya merupakan biimplikasi. Dalam
perumusan baku aturan $\LeftR{\land}$ dan~$\RightR{\lor}$, implikasi yang
bersesuaian bukan biimplikasi. Namun, terdapat versi alternatif dari aturan
tersebut yang menghasilkan biimplikasi:

\begin{defish}
    \Axiom$!A, !B, \Gamma \fCenter \Delta$
    \RightLabel{\LeftR{\land}}
    \UnaryInf$!A \land !B, \Gamma \fCenter \Delta$
    \DisplayProof
    \hfill
    \Axiom$ \Gamma \fCenter \Delta, !A, !B$
    \RightLabel{\RightR{\lor}}
    \UnaryInf$ \Gamma \fCenter \Delta, !A \lor !B$
    \DisplayProof
\end{defish}

Jika gagasan dasar ini diterapkan kepada suatu logika bernilai~$n$, hasilnya
ialah kalkulus sekuen dengan $n$ posisi, bukan dua posisi: satu posisi bagi
setiap nilai kebenaran. Bagi suatu logika bernilai tiga dengan
$V=\{\False,\Undef,\True\}$, sekuen merupakan ekspresi
$\Gamma\mid\Pi\mid\Delta$. Sekuen ini dipenuhi oleh
!!a{valuation}~$\pAssign v$ jika dan hanya jika
$\pValue{v}(!A)=\False$ bagi suatu $!A\in\Gamma$, atau
$\pValue{v}(!A)=\True$ bagi suatu $!A\in\Delta$, atau
$\pValue{v}(!A)=\Undef$ bagi suatu $!A\in\Pi$. Akibatnya, sekuen awal
$!A\mid !A\mid !A$ selalu dipenuhi.

\end{document}
